
\section*{D2 – Normal Mode Expansion and Dispersion}
For linearized tension potential \( V \sim T (\partial_\lambda \xi)^2 \), the normal mode decomposition yields:
\[
\xi^\mu(\lambda) = \sum_k a_k^\mu e^{ik\lambda}
\]
With dispersion:
\[
\omega_k \propto k
\]


\section*{Derivation D2: Normal Mode Expansion and Dispersion Relation of O1}
\subsection*{1. Linearization of the Tension Potential}
Expanding \(V_{\rm tension}[\xi]\) to second order around the equilibrium \(\xi=0\),
\[
V_{\rm tension}[\xi] \approx \frac{1}{2} \int d\lambda\,d\lambda'\,\xi^\mu(\lambda)\,K_{\mu\nu}(\lambda,\lambda')\,\xi^\nu(\lambda'),
\]
where the quadratic kernel \(K_{\mu\nu}(\lambda,\lambda')\) defines the linearized tension operator.
In many cases (e.g.\ for uniform tension \(T\) and periodic boundary conditions on a filament of length \(L\)),
\[
K_{\mu\nu}(\lambda,\lambda') = -T\,\delta_{\mu\nu}\,\frac{\partial^2}{\partial\lambda^2}\,\delta(\lambda-\lambda').
\]

\subsection*{2. Eigenvalue Problem}
The normal modes \(\phi_n^\mu(\lambda)\) satisfy
\[
\int_0^L d\lambda'\,K_{\mu\nu}(\lambda,\lambda')\,\phi_n^\nu(\lambda')
= m\,\omega_n^2\,\phi_n^\mu(\lambda).
\]
For periodic boundary conditions, plane-wave solutions
\(\phi_n^\mu(\lambda)=\frac{1}{\sqrt{L}}\,\varepsilon^\mu\,e^{ik_n\lambda}\),
with \(k_n=\frac{2\pi n}{L}\), yield the dispersion relation
\[
\omega_n = |k_n|\,\sqrt{\frac{T}{m}}.
\]

\subsection*{3. Mode Expansion of Field Operators}
Using the orthonormality of modes, the operators are expanded as
\[
\hat{\xi}^\mu(\lambda) = \sum_n \frac{1}{\sqrt{2m\omega_n}}\Bigl(\hat{a}_n\,\phi_n^\mu(\lambda)
+\hat{a}_n^\dagger\,\phi_n^{\mu*}(\lambda)\Bigr),
\]
\[
\hat{\pi}^\mu(\lambda) = -i\,\sum_n \sqrt{\frac{m\omega_n}{2}}\Bigl(\hat{a}_n\,\phi_n^\mu(\lambda)
-\hat{a}_n^\dagger\,\phi_n^{\mu*}(\lambda)\Bigr).
\]

\subsection*{4. Commutation Relations}
Imposing the canonical commutator
\(\bigl[\hat{\xi}^\mu(\lambda),\hat{\pi}_\nu(\lambda')\bigr]
= i\hbar\,\delta^\mu{}_\nu\,\delta(\lambda-\lambda')\)
implies
\[
[\hat{a}_n,\hat{a}_m^\dagger] = \delta_{nm},\quad
[\hat{a}_n,\hat{a}_m] = 0,\quad
[\hat{a}_n^\dagger,\hat{a}_m^\dagger] = 0.
\]
\section*{Derivation D3: Topological Mass Quantization of Composite Filament Bundles}
\subsection*{1. Topological Invariant \(Q\)}
Define
\[
Q = L_{\rm wind} + L_{\rm link} + W_{\rm writhe},
\]
normalized such that for basic structures:
\[
Q_{\rm neutrino}=1,\quad
Q_{\rm meson}=2,\quad
Q_{\rm baryon}=3.
\]
Here \(L_{\rm link}\) is the pairwise linking number, and \(W_{\rm writhe}\) the self-writhe of a closed contour :contentReference[oaicite:0]{index=0}.

\subsection*{2. Mass Suppression Formula}
The effective 3D inertial mass is suppressed by \(Q\):
\[
m_{\rm eff} = \frac{m_0}{Q},
\]
with \(m_0 \sim \frac{T}{\ell_f}\) the bare vibrational mass scale of a single filament :contentReference[oaicite:1]{index=1}.

\subsection*{3. Quantization of \(Q\) via Knot Invariants}
For canonical topologies:
\begin{itemize}
  \item \textbf{Unbound filament (n=1):} \(Q=1\).
  \item \textbf{Hopf link (n=2):} linking number \(=1\), minimal writhe \(=0\) \(\Rightarrow Q=2\).
  \item \textbf{Borromean link (n=3):} triple-link invariants sum to 3 \(\Rightarrow Q=3\) :contentReference[oaicite:2]{index=2}.
\end{itemize}
Thus
\[
m_{\rm neutrino}\sim m_0,\quad
m_{\rm meson}\sim \frac{m_0}{2},\quad
m_{\rm baryon}\sim \frac{m_0}{3}.
\]

\subsection*{4. Particle Family Mass Steps}
Correlate with observed families:
\[
\frac{m_{\rm meson}}{m_{\rm neutrino}}=\frac12,\quad
\frac{m_{\rm baryon}}{m_{\rm meson}}=\frac23.
\]
Deviations from these ideal ratios point to higher-order geometric corrections (e.g.\ shape factor \(\alpha_{\rm shape}\)).

\subsection*{5. Corrections and Extensions}
Including additional invariants like self-writhe \(W\) and mutual twist \(\tau\):
\[
Q = L_{\rm wind} + L_{\rm link} + W + \tau,
\]
leads to small mass renormalizations:
\[
m_{\rm eff}\approx \frac{m_0}{Q}\Bigl(1 + \frac{\alpha_{\rm_shape}}{Q}\Bigr).
\]

